%% -*- coding: utf-8 -*-
\documentclass[12pt,pagesize,paper=192mm:108mm,landscape]{scrbook} 
%1920x1080 1280x720
\areaset[current]{192mm}{108mm}
\usepackage{calc}
\usepackage[T2A]{fontenc}
\usepackage[utf8]{inputenc}
\usepackage[english,russian]{babel}
\usepackage{microtype}
\usepackage{misccorr}
\usepackage{cmap}
%\usepackage[unicode=true]{hyperref}
\usepackage{graphicx}
\usepackage{amssymb}
\usepackage{amsmath}
%\usepackage{srcltx}
\usepackage{textcomp}
\usepackage{xspace}
%научные символы и смайлики \smiley \frownie
\usepackage{wasysym}
\usepackage{ccicons}
\begin{document}
\begin{titlepage}
  \vspace*{-0.5em}
  \begin{center}    
    \hspace*{3em}
    \begin{minipage}[t]{3em}
      \includegraphics[width=\textwidth]{../BINP-logo}
    \end{minipage}\hfill
    \begin{minipage}{0.230\linewidth}
    \includegraphics[width=\textwidth]{../NSU-logo}
    \end{minipage}
    \hfill
    \hspace*{4.5em}

    Кафедра теоретической физики физического факультета НГУ

    \Large
    Профессор Грабовский А.\,В.

    \huge
    \textbf{Общая теория относительности}

    \Large
    Семинар № 14
    \vfill

    \normalsize
    \begin{minipage}{0.85\linewidth}
      Нормаль для гиперповерхностей $t =\text{const}$, $r =\text{const}$. Скаляр
      внешней кривизны для этих поверхностей в плоском
      пространстве. Вариация (производная) определителя метрического
      тензора. Демонстрация тензорных свойств абсолютно
      антисимметричного тензора 4-го ранга. Вычисление элемента
      площади гиперповерхности для гиперповерхностей $t=\text{const}$,
      $r=\text{const}$, для гиперповерхности общего вида
      $d\Sigma_i = \varepsilon_{ijkl}e_1^ie_2^je_3^k d^3y
      = \epsilon n_i\left|q\right|^\frac{1}{2} d^3y$.
    \end{minipage}
    \vfill

    \normalsize \ccbysa\hspace{0.5em}  Новосибирск 2022
  \end{center}
\end{titlepage}
\end{document}
